\section{Charge bounds and rotation averaging}
\label{sec:charge}

The history estimate of Theorem~\ref{thm:flags-local-readout} controls a
local observable by weighted flags.  To treat arbitrary hole number by
rotation averaging, we compare these weights with a physical one-body
observable, first with a slowly decreasing weight and then with a
summable weight.

Increase the fixed stopping flux in the history construction, if
necessary, so that it is at least eight and all pinning comparisons
above it are valid, including shifts by at most three.  Write
$q_{\rm s}=q_{\mathrm{stop}}$ and let $\mathcal R_{n,q}$ denote
the history unitary in \eqref{eq:flags-history-map}.  For $L\ge2$ and
$0<p<2$, put
\begin{equation}\label{eq:charge-definitions}
 \begin{gathered}
 f(x)=f_{L,p}(x)=\left(\frac{L}{L+x}\right)^p,\qquad
 C_q(f)=\frac{q+3}{3(q+1)}\sum_{j=0}^q f(j),\\
 S_{L,p}(\omega)=\sum_{b\in\mathcal B(\omega)}f(b).
 \end{gathered}
\end{equation}
When $q$ is divisible by three, $C_q(f)$ is the filled-state expectation
of $D_q(f)$: rotational invariance gives occupation
$(q+3)/(3(q+1))$ in every orbital.  The formula defines the scalar
$C_q(f)$ also at the other fluxes.  The operator $S_{L,p}$ is diagonal
in histories and scalar on each terminal summand.  Its sum is over the
flag multiset $\mathcal B(\omega)$, including all terminal flags at
position zero; each history has exactly $h=q+3-3n$ flags.  All operator
inequalities involving $S_{L,p}$ below are in history coordinates.  In particular, they make no measurement
assumption about flags in a zero-mode vector.

\begin{proposition}[Pooled charge]\label{prop:charge-pooled}
For each of $p=1/8$ and $p=3/2$ there are fixed $L\ge2$ and $c>0$ such
that, for every flux and every nonzero zero-mode space $Z_{n,q}$,
\begin{equation}\label{eq:charge-pooled}
 \mathcal R_{n,q}\bigl(C_q(f)\mathbf1-P_{n,q}D_q(f)P_{n,q}\bigr)\mathcal R_{n,q}^*
 \ \ge\ c S_{L,p}.
\end{equation}
Here compression to $Z_{n,q}$ is understood, and $P_{n,q}$ denotes its
orthogonal projection.  The constants are independent of $q,n$ and of
the readout axis.  When $h=0$, both sides vanish.
\end{proposition}

The obstruction to this estimate occurs at electron steps.  Along a
history, each F step contributes approximately $-f(q)/3$ to
$D_q(f)-C_q(f)\mathbf1$, supplying a credit for its new flag.  Each E
step contributes a positive scalar correction.  When at least two
flags remain, divide this scalar cost between them.  The scalar tail
estimate below leaves the strictly positive margin
\[
 \frac13-\frac{1+p}{6(1-p)}=\frac5{42}\qquad(p=1/8).
\]
With one remaining flag, there is no such margin.  We instead combine
the scalar with the transported observable and compare it to a filled
step; the one-hole comparison of Corollary~\ref{cor:flags-one-hole}
makes the resulting error summably small.  These two cases prove the
first charge estimate.

Rotation averaging then bounds integrated local-expectation differences
between two zero-mode states by $C(q+1)^{7/8}$ times the sum of their
hole numbers.  This still grows with the
area, but it is enough to prove the charge estimate again with $p=3/2$.
For a polar annulus $1+m\asymp M$, use this averaged bound while the
flux is moderate relative to $M$; the transport interaction has a
factor $(q+1)^{-2}$ that pays for the area dependence.  At larger
flux, apply the original fixed-axis history estimate directly to the
annulus.  The two resulting sums are small when the new weight scale
$L$ is large.  The summable weight then gives the area-independent
local-expectation bound at the end of this section.

\subsection{Scalar telescoping and summable error bounds}

We estimate $D_q(f)-C_q(f)\mathbf1$ from above.  An F step contributes
the scalar $C_{q-1}(f)-C_q(f)$, and an E step contributes
\begin{equation}\label{eq:charge-electron-scalar}
 \ell_q=C_{q-3}(f)+f(q)-C_q(f).
\end{equation}
For the continuous comparison define
\begin{equation}\label{eq:charge-continuous-scalar}
 F(x)=\frac1x\int_0^x f(t)\,dt,\qquad
 \ell(x)=-2F'(x)+f'(x)\quad(x>0),
\end{equation}
with the continuous extensions at zero.

\begin{lemma}[Scalar estimates]\label{lem:charge-scalars}
Fix $p\in(0,2)$ and a lower stopping flux $q_{\rm s}\ge8$.  There is a
function $\varepsilon_p(L)\to0$ as $L\to\infty$ with the following
properties, uniformly for $q\ge q_{\rm s}$:
\begin{align}
 \left|C_{q-1}(f)-C_q(f)+\frac13f(q)\right|
 &\le\varepsilon_p(L)f(q),\label{eq:charge-F-scalar}\\
 \sum_{q\ge q_{\rm s}}
 \frac{|\ell_q-\ell_{3\lfloor q/3\rfloor}|}{f(q)}
 &\le\varepsilon_p(L).\label{eq:charge-calibration-scalar}
\end{align}
The function $\ell$ is nonnegative.  If $0<p<1$, then for every
collection $\mathcal E$ of E-step fluxes along a history, and every
flag position $k$ below these steps,
\begin{equation}\label{eq:charge-scalar-tail}
 \sum_{\substack{q\in\mathcal E\\q\ge k}}(\ell_q)_+
 \le \left(\frac{1+p}{3(1-p)}+\varepsilon_p(L)\right)f(k).
\end{equation}
Changing $k$ by a fixed bounded amount, or taking $k$ in the terminal
range, only changes $\varepsilon_p(L)$.
\end{lemma}

\begin{proof}
The exact formula
\begin{equation}\label{eq:charge-scalar-exact}
 \ell_q=f(q)-\left(1+\frac2{q+1}\right)
 \frac{f(q)+f(q-1)+f(q-2)}3
 +\frac{2\sum_{j=0}^{q-3}f(j)}{(q+1)(q-2)}
\end{equation}
is obtained by separating the last three summands in $C_q$.
A useful positive form of this identity is obtained by putting
$d_j=f(j)-2f(j+1)+f(j+2)$ and $Q_q=(q+1)(q-2)$:
\begin{equation}\label{eq:charge-positive-scalar}
 \ell_q=\frac1{Q_q}\sum_{j=0}^{q-3}(j+1)(j+2)d_j
       +\frac{2q}{3(q+1)}d_{q-2}.
\end{equation}
Expanding the second differences verifies the formula coefficient by
coefficient.  In particular, $\ell_q\ge0$, and constant and affine
functions contribute zero.

Here are precise discrete-to-continuous error bounds:
\begin{align}
 |\ell_q-\ell(q)|+
 \sup_{x\in[q-3,q]}|\ell(x)-\ell(q)|
 &\le C(L+q)^{-2},\label{eq:charge-scalar-errors}\\
 |\ell_q-\ell_{q-r}|
 &\le
 \begin{cases}
 C L^{-2},&q\le L,\\
 C(F(q)+f(q))q^{-2},&q>L,
 \end{cases}\label{eq:charge-scalar-differences}
\end{align}
for $r=0,1,2$.  We give details because the sharper second estimate is
needed for the summable weight.  The identity
$d_j=\int (1-|t-j-1|)_+f''(t)\,dt$ expresses
\eqref{eq:charge-positive-scalar} as $\int_0^q \kappa_q(t)f''(t)\,dt$.
The function $\kappa_q$ is linear between consecutive integers and has values
\[
 \kappa_q(i)=\frac{i(i+1)}{Q_q}\ (0\le i\le q-2),\qquad
 \kappa_q(q-1)=\frac{2q}{3(q+1)},\qquad \kappa_q(q)=0.
\]
On $[0,q-2]$ its difference from $t^2/q^2$ is at most
$C(t/q^2+t^2/q^3)$, and on $[q-2,q]$ that difference is bounded by a
constant.  The error in its integral against $f''$ is consequently at
most
\[
 C\left(q^{-2}\int_0^q tf''(t)\,dt+
 q^{-3}\int_0^q t^2f''(t)\,dt+
 \int_{q-2}^q f''(t)\,dt\right).
\]
For $q\le L$ use $f''\le C/L^2$; for $q>L$ use
$\int_0^q tf''\le1$, $\int_0^q t^2f''\le2qF(q)$, and
$f''(t)\le C f(t)(L+t)^{-2}$.  This proves the first term of
\eqref{eq:charge-scalar-errors}.  The continuous identity
$\ell(x)=x^{-2}\int_0^x t^2f''(t)\,dt$ implies
$\ell'(x)=f''(x)-2\ell(x)/x$, which proves the oscillation estimate.

For \eqref{eq:charge-scalar-differences}, difference the positive
formula directly:
\[
\begin{split}
 \ell_q-\ell_{q-1}
 ={}&-\frac{2(q-1)}{Q_qQ_{q-1}}
      \sum_{j=0}^{q-4}(j+1)(j+2)d_j\\
 &+\frac{(q-1)(q-2)}{3q(q+1)}d_{q-3}
  +\frac{2q}{3(q+1)}d_{q-2}.
\end{split}
\]
The first sum is bounded by $CqF(q)$, as follows by its integral
representation and $\int_0^q t^2f''\le2qF(q)$.
This gives $CF(q)/q^2$ when $q>L$; the last two terms cost
$Cf(q)/q^2$.  Below $L$ use $d_j\le C/L^2$ in the displayed
formula.  A bounded number of differences proves the claimed bound.

The exact F-step identity is
\[
 C_{q-1}(f)-C_q(f)+\frac13 f(q)
 =\frac{2}{3q(q+1)}\sum_{j=0}^{q-1}\bigl(f(j)-f(q)\bigr).
\]
For $q\le L$ its ratio to $f(q)$ is $O(L^{-1})$.  For $q>L$ use
\begin{equation}\label{eq:charge-average-weight}
 F(q)\le
 \begin{cases}
 f(q)/(1-p),&p<1,\\
 C_pL/q,&p>1.
 \end{cases}
\end{equation}
The resulting ratio is $O(L^{-1})$ also when $1<p<2$, since
$L^{1-p}q^{p-2}\le L^{-1}$ for $q\ge L$.
The case $p=1$ follows from the explicit logarithmic integral and is
not needed below.  Dividing \eqref{eq:charge-scalar-differences} by
$f(q)$ and summing gives $O(L^{-1})$ below $L$.  Above $L$ it
gives $O(L^{-1})$: the summands are $O(q^{-2})$ for $p<1$ and
$O(q^{-2}+L^{1-p}q^{p-3})$ for $1<p<2$.

Twice integrating $f''$ proves
\[
 \ell(x)=\frac1{x^2}\int_0^x t^2f''(t)\,dt\ge0,
 \qquad
 \int_k^\infty\ell(x)\,dx=2F(k)-f(k).
\]
The E-step intervals $[q-3,q]$ have disjoint interiors along a history.
Consequently their continuous contribution is at most
$\frac13\int_{(k-3)_+}^{\infty}\ell(x)\,dx$.
For $p<1$, \eqref{eq:charge-average-weight} bounds this by
$\frac{1+p}{3(1-p)}f(k)$, up to $O(L^{-1})f(k)$ from the bounded
endpoint shift.  Finally, the tail sum of
\eqref{eq:charge-scalar-errors} is bounded by $C/(L+k)$, and
\[
 \frac1{L+k}\le\frac1L f(k)\qquad(0<p<1).
\]
Thus it is $O(L^{-1})f(k)$ uniformly in $k$.
This proves \eqref{eq:charge-scalar-tail}.
\end{proof}

We next list the analytic error sums that will be used in both charge
arguments.  Retain the size function $A_q$ from
Proposition~\ref{pinning:weighted-transport}.  With $u=1+m$,
$v=1+q-m$, and $w=\sqrt{uv/(q+1)}$, it satisfies
\begin{equation}\label{eq:charge-size}
 \begin{gathered}
 A_q(m)=\frac{u^2f''(m)}{(q+1)^2}
 \asymp\frac{f(m)}{(q+1)^2}\left(\frac{u}{L+u}\right)^2,\\
 B_q=\int_0^q A_q(m)\,dm,\qquad
 J_q=\int_0^q\frac{A_q(m)}v\,dm.
 \end{gathered}
\end{equation}
The longitude integral has bounded mass and is included in constants.
Proposition~\ref{pinning:weighted-transport} gives this size for each
transported change in $D(f)$.

\begin{lemma}[Uniform error sums]\label{lem:charge-sums}
For either $p=1/8$ or $p=3/2$, as $L\to\infty$,
\begin{equation}\label{eq:charge-common-sums}
 \sup_{q\ge q_{\rm s}}\frac{B_q}{f(q)}=O(L^{-1}),\qquad
 \sum_{q\ge q_{\rm s}}
 \left(\frac{J_q}{f(q)}+\frac1{(L+q)^2}\right)
 =O\left(\frac{1+\log L}{L}\right).
\end{equation}
For $p=1/8$ there are, in addition,
\begin{equation}\label{eq:charge-first-sums}
 \sum_{q\ge q_{\rm s}}
 \frac{\sup_{0\le m\le q}A_q(m)}{f(q)}=O(L^{-1}),\qquad
 \sum_{q\ge q_{\rm s}}\frac1{f(q)}
 \int_0^q\frac{A_q(m)}w\,dm=O(L^{-1/2}).
\end{equation}
All implied constants are independent of $L$.
\end{lemma}

\begin{proof}
For $q\le L$, integrating $u^2/L^2$ gives
$B_q\le C(q+1)L^{-2}$; for $q>L$, dropping the factor
$(u/(L+u))^2$ gives $B_q\le C F(q)/q$.
These estimates and \eqref{eq:charge-average-weight} prove the first
bound in \eqref{eq:charge-common-sums}.
On the northern half of the sphere $v\ge(q+1)/2$, whereas on the
southern half $A_q(m)\le C f(q)(L+q)^{-2}$.  Therefore
\[
 \frac{J_q}{f(q)}\le
 \frac{C B_q}{(q+1)f(q)}+
 \frac{C\log(q+2)}{(L+q)^2}.
\]
The first term sums to $O(L^{-1})$: above $L$ it is bounded by
$Cq^{-2}$ for $p<1$ and by $CL^{1-p}q^{p-3}$ for $p>1$.
The second sums to $O((1+\log L)/L)$.

For the supremum in \eqref{eq:charge-first-sums}, use $CL^{-2}$
when $q\le L$ and $C(q+1)^{-2}$ otherwise.  The latter divided by
$f(q)$ is $O(L^{-p}q^{p-2})$, whose sum is $O(L^{-1})$ for $p<1$.
On the northern half, $w\ge c\sqrt u$, so the last integral is at
most
\[
 \frac{C}{(q+1)^2}\int_0^{q/2}
 \frac{f(m)u^{3/2}}{(L+u)^2}\,dm
 +\frac{C\sqrt{q+1}\,f(q)}{(L+q)^2}.
\]
After division by $f(q)$, both terms are bounded by
$C\sqrt{q+1}/L^2$ for $q\le L$ and by $Cq^{-3/2}$ for $q>L$,
provided $p<1/2$.  Their sums are $O(L^{-1/2})$.
\end{proof}

\subsection{Calibration and the first charge estimate}

Set $K_q:=K_{q,3}(f)$ on $\mathcal F_{q-3}$.  Its compression to
$Z_{n-1,q-3}$ is the transported change in an E step from $Z_{n,q}$,
as in Proposition~\ref{prop:flags-weighted-step}.  Its center terms
have size $A_q(m)$.  A direct estimate of $\ell_q$ would accumulate a scalar
even in a filled state.  We instead calibrate $K_q+\ell_q$ by a
filled step at the nearest lower flux divisible by three.

Put $q^0=3\lfloor q/3\rfloor$ and $r=q-q^0\in\{0,1,2\}$.
Start with the filled state at flux $q^0-3$ and increase the flux
$r$ times by F transport at the same south pole.  The resulting unit
vector $\chi_q$ lies at flux $q-3$ and has the minimum possible hole
number $r$.  If $r=0$, it is the filled state itself.

\begin{lemma}[Calibrated electron step]\label{lem:charge-calibration}
For $q$ above the fixed stopping range,
\begin{equation}\label{eq:charge-calibrated-scalar}
 \left|\ell_q+\langle K_q\rangle_{\chi_q}\right|
 \le C\left(J_q+|\ell_q-\ell_{q^0}|
              +\frac{f(q)}{(L+q)^2}\right).
\end{equation}
For any collection of northern center terms of $K_q$, their expectation
in $\chi_q$ may instead be replaced by the linear reference functional
of Theorem~\ref{thm:flags-local-readout} at child flux $q-3$; the
additional error is at most $CJ_q$.  This remains true when the choice between the two centerings
is made separately on disjoint annuli.
\end{lemma}

\begin{proof}
In the filled E step at $q^0$, both parent and child compressions of
$D(f)-C(f)\mathbf1$ vanish.  Proposition~\ref{prop:flags-weighted-step}
therefore gives
$|\ell_{q^0}+\langle K_{q^0}\rangle|\le
C f(q^0)(L+q^0)^{-2}$.
Corollary~\ref{pinning:calibration-cost} compares this filled-step
expectation first with the directly embedded lower-flux state and then
with its at most two F-growth transports.  Its bound is $CJ_q$ for
each comparison, proving \eqref{eq:charge-calibrated-scalar} after
adding $|\ell_q-\ell_{q^0}|$.

For a northern source center, the linear reference functional of
Theorem~\ref{thm:flags-local-readout}, applied at child flux $q-3$, is
the expectation in the same directly embedded filled state at $q^0-3$,
using the common slots.
The last assertion of Corollary~\ref{pinning:calibration-cost} applies
to any selected collection of source pieces and bounds the sum of
absolute comparison errors by $CJ_q$.  It therefore also permits the
annulus-dependent choice of centering in the statement.
\end{proof}

For completeness, the elementary rule for summing history bounds is
as follows.  At a node of flux $q$ with positive hole number, every
descendant flag is at an index at most $q$, and hence
\begin{equation}\label{eq:charge-node-lower-bound}
 S_{L,p}\ge h f(q)\mathbf1\ge f(q)\mathbf1
\end{equation}
on that node's full descendant space.  A node error bounded by
$\eta_q f(q)\mathbf1$ can consequently be replaced by
$\eta_q S_{L,p}$ there.  If $\sum_q\eta_q\le\eta$, the sum of all
these node errors is at most $\eta S_{L,p}$: on each history, each
flag appears only at ancestor fluxes, each visited at most once.
This is a comparison of diagonal operators after the node form
bounds have been applied.  It also controls an off-diagonal
remainder by its operator norm.  At a node with no holes, use the
exact vanishing of the centered compression instead of this rule.

\begin{proof}[Proof of Proposition~\ref{prop:charge-pooled} for $p=1/8$]
Apply Proposition~\ref{prop:flags-weighted-step} recursively to
$D_q(f)-C_q(f)\mathbf1$.  The norm remainder at a node is at most
$C f(q)(L+q)^{-2}$, and its total cost is $o_L(1)S_{L,p}$ by
\eqref{eq:charge-node-lower-bound} and
Lemma~\ref{lem:charge-sums}.  F transport has norm at most $CB_q$.
Together with \eqref{eq:charge-F-scalar}, its contribution is
\[
 -\bigl(\tfrac13-o_L(1)\bigr)f(q)
\]
for the flag created at that step.

At an E step with one remaining hole, $q\equiv1\pmod3$, so $\chi_q$
also has one hole.  Compare the child vector to $\chi_q$ using Corollary~\ref{cor:flags-one-hole}.
The difference of the $K_q$ expectations is at most
$C\sup_m A_q(m)$.  This is a two-sided form bound since it holds
for every unit vector in the child space.  The calibration error
and all one-hole errors have total cost $o_L(1)S_{L,p}$ by
Lemmas~\ref{lem:charge-scalars} and \ref{lem:charge-sums}.

At an E step with at least two remaining holes, Proposition~\ref{pinning:one-sided}
bounds the increase from $K_q$ by
$C\int A_q(m)/w\,dm$.  Its total cost is again $o_L(1)S_{L,p}$.
It remains to pay the positive part of $\ell_q$.
Divide that scalar equally among the remaining flags.  A fixed flag
then pays at most one half of each $(\ell_q)_+$ at such an ancestor.
The E-step intervals are disjoint, so
Lemma~\ref{lem:charge-scalars} bounds the total payment for a flag at
$k$ by
\[
 \left(\frac{1+p}{6(1-p)}+o_L(1)\right)f(k)
 =\bigl(\tfrac3{14}+o_L(1)\bigr)f(k).
\]

At terminal fluxes there are only finitely many spaces.  Uniformly
on them, $f(j)=1+O(L^{-1})$, and hence
\[
 D_q(f)-C_q(f)\mathbf1=-\frac h3\mathbf1+O(L^{-1}).
\]
For $h=0$ the expression is exactly zero by rotational invariance;
for $h\ge1$ the error is $O(L^{-1})h$, which is also
$O(L^{-1})$ times the terminal flag weight.  Terminal flags
therefore receive the same credit $1/3-o_L(1)$.
Combining all bounds yields
\[
 \mathcal R_{n,q}(D_q(f)-C_q(f)\mathbf1)\mathcal R_{n,q}^*
 \le-\bigl(\tfrac13-\tfrac3{14}-o_L(1)\bigr)S_{L,p}.
\]
The strict margin is $5/42$.  Choose $L=L_0$ so that the total error
is less than $5/84$ and take $c=5/84$.
\end{proof}

We record what this first estimate gives after rotation averaging.
It will be applied to annular portions of the next transport
interaction, with their center size factored out.

\begin{corollary}[First averaged comparison]\label{cor:charge-weak-average}
Let $W$ be a neutral Hermitian interaction with bounded strength per
unit area and fixed rapid locality bounds.  For unit vectors
$\psi\in Z_{n,q}$ and $\chi\in Z_{n',q}$ with hole numbers $h,h'$,
\begin{equation}\label{eq:charge-weak-average}
 |\langle W\rangle_\psi-\langle W\rangle_\chi|
 \le C_W(q+1)^{7/8}(h+h').
\end{equation}
For radius-$r$ interactions of unit local norm one can take
$C_W\le C(1+r)^B$ for a fixed $B$.
\end{corollary}

\begin{proof}
To specify the rotation average on a physical zero space, let
$V_{n,q}(g)$ denote its rotation representation and define
\[
 S_{L,p}^{g}=V_{n,q}(g)\mathcal R_{n,q}^*S_{L,p}
             \mathcal R_{n,q}V_{n,q}(g)^*.
\]
Thus $S_{L,p}^{g}$ acts on $Z_{n,q}$, whereas $S_{L,p}$ itself acts
in history coordinates.  This is the history construction obtained by
rotating a fixed readout frame by $g$; its estimates are uniform in
$g$.  Use this convention separately in each number sector.

For a local term centered at $x(g)=g\cdot\mathrm{north}$, use
Theorem~\ref{thm:flags-local-readout} in the frame rotated by $g$,
with exponent $a=7/4$.  The local term itself need not belong to a
rotation-covariant family; the estimate is pointwise in $g$.  Since
$(1+b)^{-a}\le C f_{L_0,1/8}(b)$, the local comparison is bounded
by the corresponding rotated $S_{L_0,1/8}$.  Its reference scalar
cancels between the two vectors.  By irreducibility of $U_q$,
normalized rotation averaging gives the one-particle identity
\begin{equation}\label{eq:charge-rotation-average}
 \int_{SU(2)} D_q(f)^{g}\,dg
 =\frac{\sum_{j=0}^q f(j)}{q+1}\,\mathcal N,
\end{equation}
where Haar measure has mass one.  Proposition~\ref{prop:charge-pooled}
for $p=1/8$ consequently gives
\[
 \int_{SU(2)}\langle S_{L_0,1/8}^{g}\rangle_\psi\,dg
 \le\frac{h}{3c(q+1)}\sum_{j=0}^q f_{L_0,1/8}(j).
\]
The physical area is $2\pi q$, and the center density is bounded.
Multiplying by that area and using
$\sum_{j=0}^q f_{L_0,1/8}(j)\le C(q+1)^{7/8}$ proves the result.
The same argument applies to any center measure with bounded mass
in unit balls: spread each center uniformly over a ball of fixed
radius and regard its observable as centered at the new point.  The
resulting measure has bounded area density, and the support radius
increases by only a fixed constant.  Thus the preceding integral
estimate applies to both discrete and continuous center measures.  Finally,
a rapid shell expansion multiplies the radius-$r$ estimate by a summable series of shell
norms times $(1+r)^B$.
\end{proof}

\subsection{The summable weight}

We now fix $p=3/2$, keeping the already chosen $L_0$ and the constant
in Corollary~\ref{cor:charge-weak-average} fixed.  The new scale $L$
will be chosen at the end.  Calibrate every E step with positive
hole number by Lemma~\ref{lem:charge-calibration}.  Its scalar error,
the branch remainders, and all F-step errors have total cost
$o_L(1)S_{L,3/2}$ by the estimates already proved.  It remains to
bound the centered E interactions.

Divide centers into annuli $M\le1+m<2M$, where $M=1,2,4,\ldots$.
Write
\[
 a_M=\left(\frac{M}{L+M}\right)^2,\qquad f_M=f(M),\qquad
 \delta=\frac1{24}.
\]
On this annulus the local interaction size is at most
$C(q+1)^{-2}f_Ma_M$.  Choose a fixed sufficiently large $C_0$ and
separate fluxes at $T_M=C_0M^{1+\delta}$.
An annulus is present only if $q+1\ge M$.

\begin{proof}[Proof of Proposition~\ref{prop:charge-pooled} for $p=3/2$]
First consider $q\le T_M$.  Center the annular interaction at
$\chi_q$.  If the child has $h>0$ holes, the calibration state has
$r\le h$ holes, so Corollary~\ref{cor:charge-weak-average}, applied
at child flux $q-3$, bounds the centered interaction in absolute
quadratic form by
\begin{equation}\label{eq:charge-near-cost}
 C(q+1)^{-2}f_M a_M(q+1)^{7/8}h\mathbf1.
\end{equation}
The child flag weight is at least $hf(q)$.
For $M\le q+1\le T_M+1$,
$f_M/f(q)\le C M^{p\delta}$.  Thus summing the relative cost
\eqref{eq:charge-near-cost} over these fluxes gives
\begin{equation}\label{eq:charge-near-sum}
 C a_M M^{p\delta}\sum_{q+1\ge M}(q+1)^{-1-1/8}
 \le C a_M M^{-(1/8-p\delta)}
 =C a_M M^{-1/16}.
\end{equation}
The finitely many fluxes near the stopping range obey the same
estimate after increasing its constant.

Now let $q>T_M$.  With $C_0$ sufficiently large, the annulus lies
in the common northern region.  Center it by the linear reference
functional in Theorem~\ref{thm:flags-local-readout}.
Lemma~\ref{lem:charge-calibration} accounts for this change of
centering with the already budgeted error $CJ_q$.
For a center in the annulus, each descendant history $\omega$ satisfies
\begin{equation}\label{eq:charge-weight-comparison}
 \begin{split}
 W_m(\omega)
 &\le C\sum_{k\in\mathcal B(\omega)}
          \min\left(1,\left(\frac{M}{1+k}\right)^a\right)\\
 &\le\frac C{f_M}\sum_{k\in\mathcal B(\omega)}f(k)
  =\frac C{f_M}S_{L,3/2}(\omega).
 \end{split}
\end{equation}
This is also a diagonal operator inequality in the descendant history
space.
Indeed, for $k\le M$, $f(k)\ge f(M)$; for $k>M$, use $a=7/4>p$
and $M/(1+k)\le C(L+M)/(L+k)$.  The annulus has area $O(M)$.
The local readout estimate therefore bounds its centered
interaction by
\[
 C(q+1)^{-2} f_Ma_M\frac M{f_M}S_{L,3/2}.
\]
For rapid interactions this formula follows by applying the
radius-$r$ bound to each shell and summing its norm times
$(1+r)^B$.  The shell center stays in the same annulus even when
its support extends beyond it; shells reaching the south are
covered by the large-radius part of the readout estimate.
Summing over $q>T_M$ gives
\begin{equation}\label{eq:charge-far-sum}
 C a_M M\sum_{q>T_M}(q+1)^{-2}S_{L,3/2}
 \le C a_M M^{-\delta}S_{L,3/2}.
\end{equation}

Both resulting sums over dyadic $M$ tend to zero with $L$.
Explicitly, for every $0<\eta<2$,
\begin{equation}\label{eq:charge-dyadic-sum}
 \sum_{M\in\{1,2,4,\ldots\}}
 \left(\frac M{L+M}\right)^2M^{-\eta}
 \le C_\eta L^{-\eta}.
\end{equation}
For $M\le L$ this follows by summing $L^{-2}M^{2-\eta}$;
for $M>L$ it follows by summing $M^{-\eta}$.
Apply this with $\eta=1/16$ and $\eta=1/24$ to
\eqref{eq:charge-near-sum} and \eqref{eq:charge-far-sum}.
The node summation rule after
\eqref{eq:charge-node-lower-bound} now gives an error
$o_L(1)S_{L,3/2}$ for the entire centered E contribution.

There is no uncompensated electron scalar in this argument.
Every F flag and every terminal flag contributes
$-(1/3-o_L(1))f(k)$, while all E contributions and all remainders
cost $o_L(1)S_{L,3/2}$.  Choose a fixed $L=L_1$, after $L_0$ and
the constant in Corollary~\ref{cor:charge-weak-average}, so that
the combined error is at most $1/6$.  This proves
\eqref{eq:charge-pooled} with $c=1/6$.
\end{proof}

\subsection{Rotation averaging and local interactions}

The summability of $f_{L_1,3/2}$ removes the power of the area in
Corollary~\ref{cor:charge-weak-average}.  We state the resulting
estimate in the form needed for the energy argument.

\begin{theorem}[Local expectations are linear in the number of holes]
\label{thm:charge-local-linear}
Let $q=3(N-1)$ and let $W$ be a neutral Hermitian interaction whose
center density and rapid locality bounds are uniform in $q$.
For every unit vector $\psi\in Z_{n,q}$,
\begin{equation}\label{eq:charge-local-linear}
 \bigl|\langle W\rangle_\psi-
          \langle W\rangle_{\mathrm{Laughlin},N}\bigr|
 \le C_W(N-n).
\end{equation}
For interactions supported in radius $r$ with local norm at most
one, $C_W$ is bounded by $C(1+r)^B$ for a fixed exponent $B$.
The same conclusion holds for zero-mode density matrices by
linearity, with the expected value of $N-\mathcal N$ on the right.
\end{theorem}

\begin{proof}
Use the local readout estimate with $a=7/4$, now comparing
$W_0$ to $S_{L_1,3/2}$.  Since the flux is divisible by three,
the reference is the filled Laughlin expectation.
The rotation average \eqref{eq:charge-rotation-average} and the
second part of Proposition~\ref{prop:charge-pooled} give
\[
 \int_{SU(2)}\langle S_{L_1,3/2}^{g}\rangle_\psi\,dg
 \le\frac{h}{3c(q+1)}\sum_{j=0}^q f_{L_1,3/2}(j)
 \le\frac{C h}{q+1},
\]
because $L_1$ is fixed and $\sum_{j\ge0}f_{L_1,3/2}(j)<\infty$.
Integration over centers costs area $O(q+1)$, exactly as in the
proof of Corollary~\ref{cor:charge-weak-average}.  The result is
$C_Wh=3C_W(N-n)$; absorb the factor three into $C_W$.
Polynomial radius costs and rapid shell summability are preserved
by that argument.  Decomposing a density matrix into number
sectors and then into pure states proves the last assertion.
\end{proof}
